\documentclass{article}
\usepackage[utf8]{inputenc}
\usepackage{geometry}
\geometry{
  a4paper,
  total={170mm,257mm},
  left=20mm,
  top=20mm,
}
\usepackage{graphicx}
\usepackage{titling}
\usepackage{booktabs}
\usepackage{listings}
\usepackage{xcolor}
\usepackage{hyperref}
\usepackage{fancyvrb}
\usepackage{float}
\usepackage{amsmath}

\definecolor{darkblue}{rgb}{0.0, 0.0, 0.5}
\RecustomVerbatimEnvironment{verbatim}{Verbatim}{formatcom=\color{darkblue}}
\setlength{\parindent}{0pt}

\title{\textbf{DEEP LEARNING} \\
{\Large \textbf{Report 5: Backpropagation}}}
\date{May 2026}

\usepackage{fancyhdr}
\fancypagestyle{plain}{
  \fancyhf{}
  \fancyfoot[L]{\thedate}
  \fancyfoot[R]{\thepage}
  \fancyhead[L]{University of Science and Technology of Hanoi}
  \fancyhead[R]{Deep Learning}
}

\makeatletter
\def\@maketitle{%
  \newpage
  \vskip 1em
  \begin{center}
  \let \footnote \thanks
    {\LARGE \@title \par}%
    \vskip 1em
  \end{center}
  \par
  \vskip 1em}
\makeatother

\begin{document}

\pagestyle{plain}
\maketitle

\noindent\begin{tabular}{@{}ll}
  Student: & Do Minh Quang -- 2540095 \\
\end{tabular}

\section*{Implementation}

\textbf{Feedforward}

The part is done in labwork 4. This part will process inputs layer by layer, calculating the weighted sum \(z\) and activation \(a\) for each neuron.

\[
z = \sum_{i=1}^{n} x_iw_i + b
\]

All neurons use the sigmoid activation:

\[
a = \sigma(z) = \frac{1}{1 + e^{-z}}
\]


\textbf{Backpropagation}

Training relies on minimizing the Cross-Entropy loss function :
\[
J = -\frac{1}{N}\sum_{i=1}^N \left[ y_i \log(\hat{y}_i) + (1 - y_i)\log(1 - \hat{y}_i) \right]
\]

The backpropagation will compute the error in the output layer and propagate it backward to update the weights and biases of each neuron. 
The error signal \(\delta\) for each neuron is calculated based on the derivative of the loss function with respect to the neuron's output, and then used to update the parameters.
\begin{itemize}
    \item \textbf{Output layer:} $\delta = \hat{y} - y$
    \item \textbf{Hidden layers:} $\delta_j = \left(\sum_k \delta_k w_{kj}\right)\cdot \hat{y}_j(1-\hat{y}_j)$
\end{itemize}
where $w_{kj}$ is the weight from hidden neuron $j$ to neuron $k$ in the next layer, and $\delta_k$ is the error signal for neuron $k$ in the next layer.

The weights and biases are updated using the calculated error signals:
\[
    w = w - lr \times \delta \times x, \qquad b = b - lr \times \delta
\]

\textbf{Training Configuration}

The optimization of training loop use a reusable function to compute the gradient descent and the loss. 
The function will take two callbacks as arguments: one for computing the loss and another for computing the gradients.

The learning rate is set to 0.5, and the training will stop when the loss improvement is less than \(1e-6\) or after a maximum of 10000 epochs.
\section*{XOR Experiment}
The model successfully converged at epoch 4000/10000 with the detailed results as follows:

\begin{table}[H]
\centering
\caption{XOR Truth Table vs Model Predictions}
\vspace{0.2cm}
\begin{tabular}{cccc}
\toprule
\textbf{$X_1$} & \textbf{$X_2$} & \textbf{$y$} & \textbf{$\hat{y}$} \\
\midrule
0 & 0 & 0 & 0.002586 \\
0 & 1 & 1 & 0.998124 \\
1 & 0 & 1 & 0.998118 \\
1 & 1 & 0 & 0.002663 \\
\bottomrule
\end{tabular}
\end{table}

\end{document}
